
% ============================================================
% Project: The Philosophy of Viability
% Book: Book I — Ontology of Viability
% Book ID: VS.B01
% Part: Part I — Identity and Manifestation
% Part ID: VS.B01.P01
% Chapter: Chapter 08 — Comparative Philosophical Examination
% Chapter ID: VS.B01.P01.C08
% File: 2_MAINMATTER\PART_I_IDENTITY_AND_MANIFESTATION\
%       CHAPTER_08_COMPARATIVE_PHILOSOPHICAL_EXAMINATION\
%       sections\section_18_nagarjuna.tex
% Section: 18 — Nagarjuna
% Structural Identifier: VS.B01.P01.C08.S18
%
% Purpose:
% Reconstruct Nagarjuna's analysis of svabhava, dependent
% origination, emptiness, dependent designation, and the
% two truths; examine the challenge these arguments pose
% to the ontology of identity developed in Book I.
%
% Method:
% Eight-stage dialectical comparison, following Section 01
% of Chapter 08.
%
% ============================================================

\section{\texorpdfstring{N\={a}g\={a}rjuna: Dependent Origination, Emptiness, and the Problem of Identity}{Nagarjuna: Dependent Origination, Emptiness, and the Problem of Identity}}
\label{sec:ch08_s18}

% ------------------------------------------------------------
\subsection{Comparative Scope and the Shared Ontological Question}
\label{subsec:ch08_s18_ss01}
% ------------------------------------------------------------

\phantomsection\label{para:VS.B01.P01.C08.S18.SS01.P001}
The comparison with Nagarjuna concerns a fundamental difficulty in the ontological account of identity: how can an existent be determinately what it is without possessing a nature that establishes itself independently of the conditions through which it arises, exists, and becomes manifest? This question places Nagarjuna's critique of intrinsic nature in direct contact with the central concern of Part I. It does not, however, establish an equivalence between Nagarjuna's concept of emptiness and the ontology of identity developed in this book. The purpose of the comparison is to determine whether the latter can withstand the former's critique without abandoning its foundational distinctions.

\phantomsection\label{para:VS.B01.P01.C08.S18.SS01.P002}
The shared problem is not simply whether things change, whether they possess properties, or whether they can be distinguished from other things. It concerns the ontological status of the determination by which an existent is the existent that it is. If that determination is understood as an independently established essence, Nagarjuna's critique presents a direct challenge. If, however, identity is analytically distinguishable without being a separate entity, property, or self-subsistent ground, the relationship between the two positions becomes more complex and remains open to examination.

\phantomsection\label{para:VS.B01.P01.C08.S18.SS01.P003}
The comparison will therefore preserve three levels of inquiry. First, Nagarjuna's own arguments must be reconstructed without importing later interpretations into his original position. Second, the ontology of Book I must be represented according to its established principles rather than according to earlier formulations that equated identity with invariant structure or intrinsic determination. Third, the relationship between the two positions must be evaluated without assuming either compatibility or contradiction in advance. The result may be a limited convergence, a substantive disagreement, or an unresolved problem requiring further ontological clarification.

% ------------------------------------------------------------
\subsection{Historical and Textual Location}
\label{subsec:ch08_s18_ss02}
% ------------------------------------------------------------

\phantomsection\label{para:VS.B01.P01.C08.S18.SS02.P001}
Nagarjuna is conventionally situated in the early centuries of the Common Era and is regarded as a foundational figure in the development of Madhyamaka Buddhist philosophy. His principal philosophical work, the \emph{M\={u}lamadhyamakak\={a}rik\={a}}, develops a sustained examination of conditions, arising, motion, intrinsic nature, self, relations, emptiness, the two truths, and nirvana. The present section concentrates on those arguments that bear directly upon the ontological status of identity, rather than attempting a comprehensive reconstruction of the entire Buddhist philosophical tradition.

\phantomsection\label{para:VS.B01.P01.C08.S18.SS02.P002}
The principal textual references are drawn from the first, second, thirteenth, fifteenth, eighteenth, twenty-second, twenty-fourth, and twenty-fifth chapters of the \emph{M\={u}lamadhyamakak\={a}rik\={a}}. These chapters provide the relevant arguments concerning dependent arising, the impossibility of independently established intrinsic nature, the relation between a subject and its constituents, the two truths, and the status of emptiness itself. Garfield's translation and philosophical commentary provide the principal working edition for this comparison. Its interpretive character must be acknowledged: the translation is based on the Tibetan textual transmission and is not presented as a critical Sanskrit edition. The arguments are consequently cited by chapter and verse rather than by pagination alone.\footnote{See Jay L. Garfield, translator and commentator, \emph{The Fundamental Wisdom of the Middle Way: Nagarjuna's Mulamadhyamakakarika} (Oxford University Press, 1995), Preface. The edition should be treated as a significant philosophical translation and commentary, not as an unmediated or textually neutral rendering.}

% ------------------------------------------------------------
\subsection{What Nagarjuna Rejects: The Problem of Svabhava}
\label{subsec:ch08_s18_ss03}
% ------------------------------------------------------------

\phantomsection\label{para:VS.B01.P01.C08.S18.SS03.P001}
The central concept in Nagarjuna's critique is \emph{svabh\={a}va}, commonly translated as intrinsic nature, own-being, or essence. The term concerns a mode of existence in which something is established through its own nature rather than through dependence upon causes, conditions, relations, constituents, or other factors. The exact philosophical force of the term is subject to interpretation, and its English translations should not be treated as interchangeable without qualification. For the present comparison, the decisive feature is the claim of independent establishment.

\index[concept]{svabhava@svabh\={a}va}

\phantomsection\label{para:VS.B01.P01.C08.S18.SS03.P002}
Nagarjuna's objection is directed against the assumption that an existent possesses a self-established nature that makes it what it is independently of the conditions under which it arises and is intelligible. Such a nature would have to be established without depending upon anything other than itself. If it depended upon other conditions, its supposed independence would be compromised. If it did not depend upon any conditions, its arising, alteration, and causal involvement would become difficult to explain. The problem is therefore not merely the existence of properties but the claim that the ontological establishment of a thing can be secured through an independent intrinsic nature.

\phantomsection\label{para:VS.B01.P01.C08.S18.SS03.P003}
This distinction is essential when comparing Nagarjuna with the ontology of Book I. The term identity in this book cannot simply be substituted for \emph{svabh\={a}va}. Identity is not stipulated as an independently existing property or substance, nor is it defined as an invariant structure that remains untouched by all conditions. Nevertheless, the distinction cannot be settled by terminology alone. If the ontological role assigned to identity ultimately functions as a self-established essence, the Madhyamaka objection may apply even if the book avoids the vocabulary of substance or intrinsic nature. The comparison must therefore examine the actual commitments of each position.

% ------------------------------------------------------------
\subsection{The Fourfold Examination of Arising}
\label{subsec:ch08_s18_ss04}
% ------------------------------------------------------------

\phantomsection\label{para:VS.B01.P01.C08.S18.SS04.P001}
The first chapter of the \emph{M\={u}lamadhyamakak\={a}rik\={a}} opens by examining the possible modes through which arising might be established. Nagarjuna rejects arising from itself, from something wholly other, from both, or without a cause. The argument is not a simple assertion that nothing arises. It is an examination of whether arising can be coherently explained by presupposing an independently established nature in the entities involved. If something already exists through its own nature, its arising from itself becomes problematic. If it arises from something entirely other, the relation between cause and effect becomes difficult to establish. The combination of the first two possibilities does not remove their difficulties, while causeless arising would undermine the intelligibility of conditioned occurrence.\footnote{Nagarjuna, \emph{M\={u}lamadhyamakak\={a}rik\={a}}, I.1.}

\phantomsection\label{para:VS.B01.P01.C08.S18.SS04.P002}
The significance of this argument lies in its challenge to self-establishment. A thing cannot be explained as arising through an independent nature while simultaneously requiring the conditions that make its arising intelligible. The argument does not eliminate causal explanation; rather, it questions whether causation can be grounded in entities whose natures are already established independently of causal conditions. The analysis therefore introduces a central Madhyamaka constraint: the conditions of arising cannot be treated as external additions to an independently constituted existent if the existence of that existent is itself implicated in those conditions.

% ------------------------------------------------------------
\subsection{The Incompatibility of Svabhava and Conditional Arising}
\label{subsec:ch08_s18_ss05}
% ------------------------------------------------------------

\phantomsection\label{para:VS.B01.P01.C08.S18.SS05.P001}
The fifteenth chapter develops the critique by examining the relation between intrinsic nature and causes and conditions. If intrinsic nature arose through causes and conditions, it would no longer be intrinsic in the relevant sense, because its establishment would depend upon what produces it. If it did not arise through conditions, its relation to conditioned existence would remain unexplained. Nagarjuna also examines the notion of other-nature, or a nature established through something else, and argues that it cannot supply an independent foundation once the supposed intrinsic nature of the relevant terms has been undermined.\footnote{Nagarjuna, \emph{M\={u}lamadhyamakak\={a}rik\={a}}, XV.1--5.}

\phantomsection\label{para:VS.B01.P01.C08.S18.SS05.P002}
This argument presents a particularly important challenge to any ontology that treats identity as the independent essence of an existent. The issue is not whether an existent can be distinguished or whether it can be described as the same existent across different circumstances. The issue is whether the identity through which it is understood can be established independently of the conditions that constitute its existence. If identity is independent in precisely this strong sense, it approaches the structure that Nagarjuna rejects. If it is not independent, the ontology must explain how its ontological significance is preserved without making it a separate or self-subsistent entity.

% ------------------------------------------------------------
\subsection{Dependent Origination and the Mover--Motion Analysis}
\label{subsec:ch08_s18_ss06}
% ------------------------------------------------------------

\phantomsection\label{para:VS.B01.P01.C08.S18.SS06.P001}
The concept of \emph{prat\={i}tyasamutp\={a}da}, or dependent origination, provides the broader context for this critique. It concerns the conditioned arising of phenomena rather than the independent establishment of things that subsequently enter into relations. Dependence is not merely a relation added after the entities have already been constituted. The question is whether the entities involved can be understood as independently established prior to the conditions through which they arise and function. This is why the analysis of dependence has direct ontological implications.\footnote{See Nagarjuna, \emph{M\={u}lamadhyamakak\={a}rik\={a}}, I.1 and XXIV.18--19.}

\index[concept]{pratityasamutpada@prat\={i}tyasamutp\={a}da}

\phantomsection\label{para:VS.B01.P01.C08.S18.SS06.P002}
The second chapter's examination of the mover and motion offers a more specific illustration. The designation of something as a mover cannot be established independently of motion, while motion cannot be adequately characterized by simply presupposing an independently established mover. The analysis challenges the assumption that the agent and the activity can each be secured independently before their relation is considered. It therefore demonstrates how a particular ontological or conceptual determination may depend upon the conditions under which it becomes intelligible. This argument is significant, but it must remain limited to the relation it actually examines; it does not by itself establish that every form of identity is reducible to a relation.\footnote{Nagarjuna, \emph{M\={u}lamadhyamakak\={a}rik\={a}}, II.18--23.}

\phantomsection\label{para:VS.B01.P01.C08.S18.SS06.P003}
The mover--motion analysis consequently supplies a test rather than a complete theory of identity. It demonstrates the difficulty of treating a relationally specified existent as independently established before the relevant relation. It does not establish that identity itself is identical with relational determination, nor does it show that all ontological distinctions are merely conventional. The precise scope of the argument must be preserved. Its force lies in demonstrating that the independence of an entity cannot simply be assumed where its determination depends upon the conditions through which it acts or is described.

% ------------------------------------------------------------
\subsection{Emptiness and the Reality of Conventional Phenomena}
\label{subsec:ch08_s18_ss07}
% ------------------------------------------------------------

\phantomsection\label{para:VS.B01.P01.C08.S18.SS07.P001}
The concept of \emph{\'{s}\={u}nyat\={a}}, or emptiness, expresses the absence of independently established intrinsic nature. It should not be interpreted as a synonym for absolute nonexistence, annihilation, or the disappearance of all distinctions. Such interpretations would confuse the rejection of self-established nature with the denial of conditioned phenomena. Nagarjuna's argument concerns the mode in which phenomena exist, not a simple declaration that phenomena do not exist in any sense.\footnote{Nagarjuna, \emph{M\={u}lamadhyamakak\={a}rik\={a}}, XXIV.8--10, XXIV.18--19.}

\index[concept]{sunyata@\'{s}\={u}nyat\={a}}

\phantomsection\label{para:VS.B01.P01.C08.S18.SS07.P002}
The twenty-fourth chapter is particularly important because it addresses the objection that emptiness would destroy the possibility of causal efficacy, practice, and the Four Noble Truths. Nagarjuna's response depends upon the relation between conventional and ultimate truth. The denial of intrinsic nature does not eliminate the conventional conditions through which phenomena function. Indeed, the argument that all phenomena are empty is connected with the claim that they arise dependently. Emptiness therefore cannot be equated with a void in which nothing happens. It concerns the absence of independent establishment within the very field of conditioned occurrence.

% ------------------------------------------------------------
\subsection{The Analysis of Self and Aggregates}
\label{subsec:ch08_s18_ss08}
% ------------------------------------------------------------

\phantomsection\label{para:VS.B01.P01.C08.S18.SS08.P001}
The eighteenth chapter examines the relation between the self and the aggregates. If the self were identical with the aggregates, it would inherit their arising and disintegration. If it were entirely different from them, the relation through which the self is attributed their characteristics would become difficult to explain. The argument challenges the establishment of a self that exists independently of the constituents and conditions through which personal existence is understood.\footnote{Nagarjuna, \emph{M\={u}lamadhyamakak\={a}rik\={a}}, XVIII.1--2.}

\index[concept]{anatman@an\={a}tman}

\phantomsection\label{para:VS.B01.P01.C08.S18.SS08.P002}
This argument must not be expanded beyond its textual scope. It does not demonstrate that every conventional distinction between persons is meaningless, nor does it establish that all forms of identity are unreal. Its immediate target is the independently established self. The wider Madhyamaka critique extends the analysis of intrinsic nature to other phenomena, but the specific argument concerning self and aggregates cannot simply be treated as a complete theory of ontological identity. The distinction between personal selfhood and the broader question of the identity of an existent must therefore remain explicit.

% ------------------------------------------------------------
\subsection{The Two Truths}
\label{subsec:ch08_s18_ss09}
% ------------------------------------------------------------

\phantomsection\label{para:VS.B01.P01.C08.S18.SS09.P001}
Nagarjuna's account of the two truths distinguishes conventional truth from ultimate truth. The twenty-fourth chapter states that the teaching of the Buddhas is presented in dependence upon these two truths and emphasizes that failure to understand their distinction prevents an adequate understanding of the teaching. The relation is not a distinction between two unrelated subjects. Conventional truth is necessary for communicating and understanding the ultimate, while the ultimate cannot be articulated without relying upon conventional discourse.\footnote{Nagarjuna, \emph{M\={u}lamadhyamakak\={a}rik\={a}}, XXIV.8--10.}

\index[concept]{two truths}
\index[concept]{conventional truth}
\index[concept]{ultimate truth}

\phantomsection\label{para:VS.B01.P01.C08.S18.SS09.P002}
The two truths should not be treated as two independent ontological worlds or as a simple division between appearance and an underlying reality. Interpretations of their precise relation differ, and the present comparison does not require the adoption of one comprehensive account of that debate. What is essential here is that the two truths perform a function different from the analytical distinction between Being and Identity in Book I. They are not equivalent to Being and Identity, nor are they equivalent to Identity and Manifestation. Mapping these distinctions onto one another would obscure the different questions each framework addresses.

% ------------------------------------------------------------
\subsection{Dependent Designation and the Status of Identity}
\label{subsec:ch08_s18_ss10}
% ------------------------------------------------------------

\phantomsection\label{para:VS.B01.P01.C08.S18.SS10.P001}
The most important passage for the present comparison occurs in the twenty-fourth chapter, where dependent origination, emptiness, dependent designation, and the Middle Way are brought into a single argument. The expression commonly translated as dependent designation indicates that phenomena are understood through dependence rather than through independently established intrinsic nature. The passage does not authorize the simple conclusion that things are nothing more than words or that an individual act of cognition creates their existence. It places the problem of designation within the broader argument concerning dependence and emptiness.\footnote{Nagarjuna, \emph{M\={u}lamadhyamakak\={a}rik\={a}}, XXIV.18. The English expression ``dependent designation'' is an interpretive rendering of the relevant Sanskrit compound.}

\index[concept]{dependent designation}

\phantomsection\label{para:VS.B01.P01.C08.S18.SS10.P002}
Nevertheless, the passage raises a difficult question for any ontology that assigns identity an ontological status. If the determination of an existent is dependent upon the conditions through which it is designated, how should the relation between the existent and its identity be understood? Does dependent designation concern only the conventional articulation of phenomena, or does it also disclose something about the conditions of their ontological establishment? The text does not permit these questions to be answered by simply identifying designation with manifestation or by reducing dependence to a cognitive operation. The precise ontological significance of designation remains a matter requiring careful interpretation.

% ------------------------------------------------------------
\subsection{The Catuskoti and the Limits of Predication}
\label{subsec:ch08_s18_ss11}
% ------------------------------------------------------------

\phantomsection\label{para:VS.B01.P01.C08.S18.SS11.P001}
The \emph{catu\d{s}ko\d{t}i}, commonly called the tetralemma or fourfold analysis, is associated with Nagarjuna's examination of claims that phenomena possess independently established natures. Its application must, however, be treated with precision. The four modes of arising examined in the first chapter are arising from itself, from another, from both, and without cause. These are not simply identical with the four possibilities concerning whether something exists, does not exist, both exists and does not exist, or neither exists nor does not exist. The fourfold structure must be interpreted in relation to the specific question under examination.\footnote{Compare Nagarjuna, \emph{M\={u}lamadhyamakak\={a}rik\={a}}, I.1, XXII.11, and XXV.10--18.}

\index[concept]{catuskoti@catu\d{s}ko\d{t}i}

\phantomsection\label{para:VS.B01.P01.C08.S18.SS11.P002}
The importance of the fourfold analysis is not that it supplies a universal logical rule by which every proposition is simply declared false. Its philosophical role concerns the assumptions that make certain alternatives appear independently established. The analysis exposes the difficulties involved in assigning an intrinsic nature through the available modes of predication. Consequently, the tetralemma should not be reduced to a generic four-valued logic, nor should every use of fourfold negation be attributed to the same argument. The textual context determines the force of each examination.

% ------------------------------------------------------------
\subsection{The Emptiness of Emptiness}
\label{subsec:ch08_s18_ss12}
% ------------------------------------------------------------

\phantomsection\label{para:VS.B01.P01.C08.S18.SS12.P001}
A further difficulty arises if emptiness itself is treated as an independently existing nature. Such an interpretation would reproduce the very structure that the critique of intrinsic nature seeks to undermine. Nagarjuna's warning against turning emptiness into a view is therefore essential. Emptiness cannot be introduced as a new substance, ultimate entity, or self-established ontological ground that replaces the intrinsic natures previously rejected. Its own conceptual status must remain subject to the critique of independent establishment.\footnote{Nagarjuna, \emph{M\={u}lamadhyamakak\={a}rik\={a}}, XIII.8 and XXII.11.}

\phantomsection\label{para:VS.B01.P01.C08.S18.SS12.P002}
This point prevents a common misunderstanding of Madhyamaka. Emptiness is not a hidden reality behind phenomena, nor is it an additional property possessed by every existent. If it were treated in either way, the analysis would merely replace one form of reification with another. The relation between emptiness and dependent origination must therefore be understood without converting emptiness into an independent foundation. This self-limiting character is particularly important when comparing Nagarjuna with any ontology that seeks to articulate a positive account of identity without reducing it to an independently existing essence.

% ------------------------------------------------------------
\subsection{Interpretive Disputes and Textual Caution}
\label{subsec:ch08_s18_ss13}
% ------------------------------------------------------------

\phantomsection\label{para:VS.B01.P01.C08.S18.SS13.P001}
The preceding reconstruction should not be mistaken for a complete or uncontested interpretation of Nagarjuna. Several important questions remain disputed. The meaning of \emph{svabh\={a}va}, the scope of emptiness, the relationship between conventional and ultimate truth, and the precise role of the tetralemma have all received different interpretations. The philosophical significance of Nagarjuna's arguments also depends upon whether they are understood primarily as a critique of independent existence, as a dialectical method for dismantling philosophical positions, or as a more comprehensive account of dependent phenomena. These approaches overlap but should not be treated as identical.

\phantomsection\label{para:VS.B01.P01.C08.S18.SS13.P002}
The translation of the root text introduces an additional limitation. Garfield's interpretation is shaped by a particular Madhyamaka reading and by the Tibetan textual transmission. His commentary is philosophically substantial, but it cannot be treated as an unmediated statement of everything that Nagarjuna intended. The present section therefore distinguishes between the wording of a translated verse, the interpretive claims made in its commentary, and the comparative conclusions drawn from it. No conclusion about the identity of an existent will be based solely upon a single English rendering of a contested term.

\phantomsection\label{para:VS.B01.P01.C08.S18.SS13.P003}
The distinction between Nagarjuna's own arguments and the later development of Madhyamaka is equally important. Later interpretations may clarify the philosophical possibilities opened by the root text, but they cannot automatically be attributed to Nagarjuna himself. In particular, distinctions developed within later Madhyamaka traditions should not be projected backward as though they were already fully formulated in the original work. The present comparison therefore remains focused on the arguments of the \emph{M\={u}lamadhyamakak\={a}rik\={a}} and uses later philosophical interpretation only where it helps clarify an issue rather than replace the textual reconstruction.

% ------------------------------------------------------------
\subsection{Nagarjuna's Strongest Objection to the Ontology of Book I}
\label{subsec:ch08_s18_ss14}
% ------------------------------------------------------------

\phantomsection\label{para:VS.B01.P01.C08.S18.SS14.P001}
The strongest objection that can be developed from Nagarjuna's analysis concerns the ontological status assigned to identity in Book I. The book characterizes identity as what makes an existent the existent that it is, while maintaining that identity is not an additional property, an independent entity, or a determination that can be separated from Being. These safeguards are important, but they do not by themselves establish that identity is free from every form of intrinsic self-establishment. Nagarjuna could ask whether the identity principle nevertheless functions as the essence through which an existent is established as itself.

\phantomsection\label{para:VS.B01.P01.C08.S18.SS14.P002}
The objection becomes more difficult when the three foundational principles are considered together. Analytical Distinction maintains that Being and Identity are conceptually distinguishable. Ontological Inseparability denies that they exist as independently separable realities. Absence of Ontological Priority denies that either has ontological precedence over the other. Yet Nagarjuna could still question whether the identity dimension, even without being a separate substance, functions as an independently established nature of the existent. In other words, inseparability and the absence of priority do not automatically demonstrate the absence of intrinsic establishment.

\phantomsection\label{para:VS.B01.P01.C08.S18.SS14.P003}
A further question concerns the positive content of identity. If identity is not a property, not determination, not unity, not structure, not invariance, not composition, and not relational manifestation, what establishes its ontological significance? A sequence of exclusions cannot by itself provide a complete positive account. If identity is indispensable to the ontology, its relation to the conditions through which an existent arises and is constituted must be explained. Otherwise, the concept risks functioning as an unexamined ontological principle whose independence is denied verbally but retained in its explanatory role. This is the strongest form of the Madhyamaka challenge, and it cannot be dismissed by claiming that Book I does not use the word essence.

% ------------------------------------------------------------
\subsection{Book I's Strongest Response}
\label{subsec:ch08_s18_ss15}
% ------------------------------------------------------------

\phantomsection\label{para:VS.B01.P01.C08.S18.SS15.P001}
The strongest response available from Book I begins by rejecting the assumption that identity must be understood either as a self-subsistent essence or as nothing more than a conventional designation. The foundational account does not introduce identity as an additional entity that exists alongside Being. Nor does it claim that identity is an independently existing bearer of properties. The distinction between Being and Identity is analytical, while their ontological inseparability denies that either can be established as a separate reality. The absence of ontological priority further prevents either from being treated as the independent foundation of the other.

\phantomsection\label{para:VS.B01.P01.C08.S18.SS15.P002}
This response also requires preserving the distinctions developed in the earlier chapters. Identity cannot be reduced to determination, even though the determination of an existent may be relevant to its articulation. It cannot be equated with invariance, because the conditions under which something remains invariant do not exhaust the question of what makes it the existent that it is. It cannot be identified with manifestation, because the way an existent becomes manifest does not by itself constitute its identity. These distinctions prevent the ontology from treating identity as a single property or structural feature that can be substituted for every other ontological category.

\phantomsection\label{para:VS.B01.P01.C08.S18.SS15.P003}
Nevertheless, this response remains conditional. The denial that identity is an independent substance or property does not yet explain how its ontological significance is sustained through dependence. The three foundational principles establish important constraints, but the question of the relation between identity and the conditions of an existent remains open. Book I therefore has a defensible response to the direct identification of identity with \emph{svabh\={a}va}, but it has not thereby demonstrated a complete compatibility with Nagarjuna's analysis. The distinction between analytical inseparability and self-established nature must be made philosophically explicit rather than assumed to be self-evident.

% ------------------------------------------------------------
\subsection{Dependence, Relationality, and Knower-Independence}
\label{subsec:ch08_s18_ss16}
% ------------------------------------------------------------

\phantomsection\label{para:VS.B01.P01.C08.S18.SS16.P001}
The most important unresolved issue concerns the meaning of dependence. Nagarjuna's critique requires that phenomena not be understood as independently established through intrinsic nature. Book I must therefore explain how identity can be ontologically inseparable from an existent without becoming independent of the conditions through which that existent arises. This does not mean that identity must be reduced to a relation, because relational determination and identity are not identical in the ontology of Book I. The question is whether dependence can be acknowledged without making identity merely the sum or product of its relations.

\phantomsection\label{para:VS.B01.P01.C08.S18.SS16.P002}
A related problem concerns the role of designation. The use of conventional designation in Nagarjuna's analysis must not automatically be interpreted as a claim that existence depends upon an individual knower, a subjective act of cognition, or a linguistic construction. Such a reading would require additional argument. At the same time, the ontological role of convention cannot be ignored. Book I maintains that existence, identity, and distinction do not depend upon cognition or language for their existence. The precise relationship between this knower-independent commitment and Nagarjuna's account of dependent designation remains a genuine point of philosophical tension.

% ------------------------------------------------------------
\subsection{Book I's Strongest Objection to Madhyamaka}
\label{subsec:ch08_s18_ss17}
% ------------------------------------------------------------

\phantomsection\label{para:VS.B01.P01.C08.S18.SS17.P001}
The strongest objection that Book I can raise against an interpretation of Madhyamaka concerns the explanatory status of conventional identity. If an existent is understood through dependent designation and lacks intrinsic nature, what explains the non-arbitrariness of the distinctions through which it is identified? How are the conditions that make a particular existent distinguishable from another to be understood without assuming that the distinctions are merely imposed by convention? The objection is not that Nagarjuna denies conventional distinctions. It is whether the account of dependence sufficiently explains their ontological significance for a theory whose central concern is identity.

\phantomsection\label{para:VS.B01.P01.C08.S18.SS17.P002}
This question must be formulated without attributing nihilism to Nagarjuna. The conventional efficacy of phenomena is not simply denied, and the analysis of emptiness is not intended to abolish causal relations. The more precise question is whether the conventional and dependent status of identity provides a sufficient account of the ontological constitution of an existent. Book I can ask whether a framework that refuses intrinsic nature can nevertheless articulate why particular distinctions are not arbitrary and how they retain their significance beyond the circumstances of a particular description.

% ------------------------------------------------------------
\subsection{Madhyamaka's Strongest Possible Reply}
\label{subsec:ch08_s18_ss18}
% ------------------------------------------------------------

\phantomsection\label{para:VS.B01.P01.C08.S18.SS18.P001}
The strongest Madhyamaka reply is that dependence does not imply arbitrariness. A phenomenon may function through conditions, causal relations, constituents, and conventional designations without possessing an independently established intrinsic nature. The absence of \emph{svabh\={a}va} does not remove the conditions through which phenomena arise or the distinctions through which they are conventionally understood. Rather, it prevents these conditions from being transformed into a self-subsistent foundation. The conventional functioning of phenomena is therefore not necessarily compromised by emptiness; it is understood through dependence rather than through intrinsic establishment.

\phantomsection\label{para:VS.B01.P01.C08.S18.SS18.P002}
From this perspective, the demand for an independently established identity may itself express the problem that Nagarjuna's critique seeks to expose. If the conditions of an existent are sufficient to explain its conventional functioning, an additional intrinsic identity may appear unnecessary. The Madhyamaka response would therefore question whether Book I's insistence upon the ontological significance of identity introduces an explanatory commitment that goes beyond what dependent origination requires. It would ask whether the ontology has demonstrated the need for identity as a distinct ontological dimension or merely preserved a familiar conceptual expectation.

\phantomsection\label{para:VS.B01.P01.C08.S18.SS18.P003}
This response is strong because it does not deny the reality of conventional distinctions. It instead challenges the assumption that their reality requires intrinsic establishment. However, it does not automatically settle the problem posed by Book I. The latter can still ask whether the conditions of dependent designation provide a sufficiently articulated account of identity, rather than merely explaining the practical use of designations. The exchange therefore remains substantive: one side questions the need for an independently articulated identity principle, while the other questions whether dependence alone exhausts the ontological significance of identity.

% ------------------------------------------------------------
\subsection{Comparative Costs, Convergences, and Open Problems}
\label{subsec:ch08_s18_ss19}
% ------------------------------------------------------------

\phantomsection\label{para:VS.B01.P01.C08.S18.SS19.P001}
If Book I accepts the force of Nagarjuna's critique, it must clarify that identity cannot be understood as an independently established essence, even when it is not described as a substance or property. It must explain how the ontological inseparability of Being and Identity remains compatible with dependence and how the distinction between them avoids becoming a new form of self-establishment. This would require a more explicit account of the conditions through which identity is ontologically significant. It would not require abandoning the three foundational principles, but it would place their interpretation under a more demanding philosophical constraint.

\phantomsection\label{para:VS.B01.P01.C08.S18.SS19.P002}
If, conversely, Madhyamaka is interpreted as denying every form of identity or treating conventional designation as an arbitrary linguistic construction, it risks losing the force of its own account of dependent phenomena. Such a reading would not adequately explain why conventional distinctions remain effective or why causal conditions are not simply interchangeable. This, however, is not the strongest interpretation of Nagarjuna. His analysis explicitly preserves the significance of conventional truth and dependent arising. The objection therefore applies primarily to reductive interpretations of Madhyamaka rather than to the most defensible reconstruction of the root text.

\phantomsection\label{para:VS.B01.P01.C08.S18.SS19.P003}
A limited convergence can nevertheless be identified. Both positions resist the idea that the ontological status of an existent can be adequately explained by positing a separate, independently established essence. Both require careful distinctions between an existent and the conditions through which it is understood. Their disagreement concerns the positive ontological status of identity and the explanatory role assigned to it. Book I seeks to articulate identity as an analytically distinguishable but ontologically inseparable dimension of Being. Nagarjuna's analysis challenges any interpretation that would allow this dimension to function as an independently established nature. Neither position can be substituted for the other without loss of its central commitments.

% ------------------------------------------------------------
\subsection{Claim Status, Unresolved Questions, and Transition}
\label{subsec:ch08_s18_ss20}
% ------------------------------------------------------------

\phantomsection\label{para:VS.B01.P01.C08.S18.SS20.P001}
The claims developed in this comparison must be distinguished according to their functional form and epistemic status. The reconstruction of Nagarjuna's critique of intrinsic nature is an argument supported by the relevant textual passages and commentary. The conclusion that \emph{svabh\={a}va} cannot simply be equated with the concept of identity in Book I is an argued comparative distinction. The claim that the three principles of Book I are compatible with the Madhyamaka critique remains unestablished. It is therefore not adopted as a foundational commitment, axiom, definition, or settled thesis. The relation between identity, dependence, and dependent designation remains an unresolved ontological problem.

\phantomsection\label{para:VS.B01.P01.C08.S18.SS20.P002}
The principal open questions are consequently these:

\begin{enumerate}
    \item Can identity be ontologically inseparable from Being without functioning as an independently established nature?
    \item What is the ontological status of an analytical distinction that does not correspond to two independently existing entities?
    \item Does dependent origination require that identity be understood through relations, and if so, how can identity remain non-reducible to relational determination?
    \item How can dependent designation be understood without reducing existence to an individual knower or a linguistic construction?
    \item Can the conventional efficacy of distinctions be explained without reintroducing intrinsic nature?
    \item How should identity be understood in relation to persistence and continuity without equating identity with either?
\end{enumerate}

\phantomsection\label{para:VS.B01.P01.C08.S18.SS20.P003}
The comparison with Nagarjuna therefore produces neither a complete reconciliation nor a decisive refutation of Book I. Its principal result is a constraint: the ontology of identity must demonstrate that ontological inseparability does not imply intrinsic self-establishment. At the same time, the critique of intrinsic nature must not be simplified into the denial of conventional identity or the reduction of existence to language and cognition. The next comparison, with Vasubandhu, will examine a different set of questions concerning the relation between phenomena, appearance, and cognition. That inquiry must be reconstructed independently and must not be prejudged as either a continuation of Nagarjuna's position or its direct negation.

% ============================================================
% END OF SECTION 18 OF CHAPTER 08
% ============================================================
